\section{Memoria activa y tiempo extramental discreto}

\subsection{Memoria como retención activa}

El trapping estructural y la hiper-persistencia no son formas de retención pasiva de estados pasados, sino las expresiones de lo que se denomina «memoria activa»:

\begin{itemize}
    \item El \textit{trapping} (LAM y TT) mide la duración de la permanencia en patrones recurrentes. Representa la capacidad de un módulo de \textit{sostener} una configuración el tiempo suficiente como para que dicha configuración pueda tener consecuencias estructurales.
    \item La \textit{hiper-persistencia} añade un componente de auto-referencia: el módulo sostiene su configuración \textbf{por encima} de su propio estándar reciente. No se trata solo de retener, sino de retener con mayor intensidad de la que el módulo exhibe habitualmente.
\end{itemize}

Esta memoria es \textbf{dinámica y estabilizadora}. No se limita a conservar el pasado, sino que mantiene activamente una configuración con la firmeza necesaria para que adquiera peso ontológico en las transiciones entre capas.

\subsection{El tiempo extramental como tiempo estratificado}

La visión del presente como fenómeno estratificado y fragmentado tiene consecuencias directas para la comprensión del tiempo extramental:

\begin{enumerate}
    \item El tiempo extramental no es un flujo uniforme ni una sucesión de instantes discretos vacíos. Es un \textbf{«tiempo estratificado por persistencia»}: un tiempo articulado por cambios en la densidad, la adhesión y la coherencia de las configuraciones de persistencia.
    
    \item La \textbf{discreción} que detecta el RECD corresponde, en lo fundamental, a \textbf{transiciones entre capas} de diferente nivel de integración, y no a instantes puntuales aislados.
    
    \item La antisincronización introduce una dimensión \textbf{relacional} en el tiempo extramental: el grado de fragmentación o de coherencia entre los presentes de los módulos condiciona la posibilidad de que ocurran articulaciones discretas significativas a nivel global.
    
    \item El tiempo extramental aparece como un \textbf{tiempo jerárquico}, en el que diferentes capas de densidad de persistencia pueden coexistir y resistirse mutuamente. Las transiciones entre estas capas constituyen los eventos ontológicamente primarios.
\end{enumerate}

Esta concepción se inscribe en una tradición que, desde Aristóteles hasta Whitehead y Prigogine, ha insistido en que el tiempo no es independiente de la dinámica de los procesos. Añade, sin embargo, la tesis de que dicha dinámica está organizada en capas jerárquicas de persistencia cuya integración está regulada por la antisincronización \citep{padilla_polo_dialogue_2026}.